\begin{solution}
Einstein--de Sitter has \(a\propto t^{2/3}\), hence
\[
t_0=\frac{2}{3H_0},\qquad
t(z)=\frac{2}{3H_0}(1+z)^{-3/2},
\]
and
\[
t_{\rm lookback}=\frac{2}{3H_0}
\left[1-(1+z)^{-3/2}\right].
\]
The age is two thirds of the Hubble time.
\end{solution}

\begin{solution}
For Einstein--de Sitter,
\[
d_A(z)=\frac{2}{H_0}\frac{1-(1+z)^{-1/2}}{1+z}.
\]
Differentiation gives the maximum at
\[
1+z=\frac94,\qquad z=\frac54.
\]
Beyond it the scale factor at emission shrinks faster than comoving distance
grows, so a fixed proper ruler subtends a larger angle.
\end{solution}

\begin{solution}
The inferred sizes are
\[
\Delta r_\parallel=\frac{\Delta z}{H(z)},\qquad
\Delta r_\perp=D_M(z)\Delta\theta.
\]
Sphericity requires their ratio to be unity, so the observable combination is
\[
F_{\rm AP}(z)=H(z)D_M(z)
=\frac{\Delta z}{\Delta\theta}
\]
for a spherical feature, with factors of \(c\) restored as needed.
\end{solution}

\begin{solution}
The expansion is
\[
d_L=\frac1{H_0}\left[
z+\frac12(1-q_0)z^2
-\frac16(1-q_0-3q_0^2+j_0+\Omega_{k0})z^3+\cdots
\right]
\]
for the stated curvature convention.  Flat \(\Lambda\)CDM has \(j_0=1\);
the cubic coefficient first probes jerk together with curvature.
\end{solution}

\begin{solution}
Using the transverse Poisson equation and geodesic deviation,
\[
\kappa(\hat{\bm n})
=\frac32H_0^2\Omega_{m0}\int_0^{\chi_s}\dd\chi\,
\frac{\chi(\chi_s-\chi)}{\chi_s}
\frac{\delta(\chi\hat{\bm n},z)}{a(\chi)}
\]
in a flat universe.  The factor
\(\chi(\chi_s-\chi)/\chi_s\) is the geometric lensing kernel.
\end{solution}

\begin{solution}
The comoving volume element gives
\[
\frac{\dd N}{\dd z\,\dd\Omega}
=n_c(z)S(z)\frac{D_M^2(z)}{H(z)}.
\]
The last factor is pure geometry, \(n_c(z)\) describes population evolution,
and \(S(z)\) contains flux limits, completeness, and other observational
selection.
\end{solution}

\begin{solution}
The spectral relation is
\[
F_{\nu_o}=\frac{1+z}{4\pi d_L^2}
L_{\nu_e}\big|_{\nu_e=(1+z)\nu_o}.
\]
For \(L_\nu\propto\nu^{-\alpha}\),
\[
F_{\nu_o}\propto\frac{(1+z)^{1-\alpha}}{d_L^2}L_{\nu_o}.
\]
Relative to the bolometric inverse-square convention, the magnitude
correction is \(K=-2.5(1-\alpha)\log_{10}(1+z)\), subject to filter
conventions.
\end{solution}

\begin{solution}
Since \(1+z=a_0/a(t)\),
\[
\frac{\dd z}{\dd t}=-(1+z)H(z),
\]
which yields the formula.  The galaxies must share a common formation clock,
evolve passively, have reliably modeled differential ages, and not have age
trends from changing metallicity or selection masquerading as cosmic time.
\end{solution}

\begin{solution}
The transverse angle and radial redshift interval are
\[
\Delta\theta=\frac{r_s}{D_M(z)},\qquad
\Delta z=H(z)r_s.
\]
Thus angular clustering measures \(D_M/r_s\), radial clustering measures
\(Hr_s\), and an isotropic analysis mainly constrains their volume-averaged
combination.
\end{solution}

\begin{solution}
To first order,
\[
z_{\rm obs}\simeq H_0d+v_\parallel,
\qquad
\frac{\delta \dd}{\dd}\simeq\frac{v_\parallel}{H_0\dd}
\simeq\frac{v_\parallel}{z}.
\]
The distance-modulus perturbation is
\(\delta\mu=(5/\ln10)\delta d/d\).  Because the velocity term is nearly
distance independent while the Hubble redshift tends to zero, its fractional
effect grows as \(1/z\).
\end{solution}
